\scanpage{361}
% @ head of sheet
\mspage{85}
% @ top, lines 1–7
que les $\Omega^n(X,x)$ (pour $n\geq1$) sont $0$-connexes. (Cf. $\overline{\underline{\mathrm{XII}}}$ cor. p. 190 déjà cité.) Pour ceci, à cause de $\mathrm{Loc}(5)$ il suffit de prouver qu'ils sont $W$-asphériques, et par induction il suffit de prouver que $\Omega^1(X,x)$ l'est. Or par définition de $\Omega^1(X,x)$, c'est là la fibre en $x$ d'une flèche dans $\mathrm{Cat}$
% @ upper middle, path-category arrow
\diagram{\begin{tikzcd}[ampersand replacement=\&]
E=\underline{\mathrm{Ch}}(X,x,-) \arrow[r,"p\in\mathrm{Parf}_{W_{\unsure{\omega|W|w}}}"] \& X{,}
\end{tikzcd}}
% @ middle, lines 9–14
où la source $E$ est $W_\omega$-asphérique, a fortiori $W$-asphérique. Comme $X$ l'est également, on a donc $p\in W$, donc $p\in\mathrm{Parf}_{W_\omega}\cap W\subset W^{\unsure{u|n}}$, donc $p$ est à fibres $W$-asphériques, et en particulier $\Omega^1(X,x)$ est $W$-asphérique, qed.

% @ lower third, Corollary 6.3
\uline{Corollaire} 6.3. (Caractérisation de $\mathbb W_\infty$). $\mathbb W_\infty$ est le plus grand localiseur fondamental satisfaisant à $\mathrm{Loc}(5,7\,D_{W_\omega})$ - donc aussi le plus grand loc. fond. satisfaisant à $\mathrm{Loc}(5,7)$.

% @ foot, final three lines
Je pense que nous trouverons plus bas une caractérisation ``interne'' de $\mathbb W_\infty$, via un ``axiome de descente'' (\deleted{\unsure{b}} \unsure{très} fort) $\mathrm{Loc}(10)$,
% archive sheet number: 360
